\documentclass[12pt,a4paper]{article}
\usepackage{fontspec}
\setmainfont{Liberation Serif}
\usepackage[english,russian]{babel}
\usepackage{amsmath,amssymb,amsthm,mathtools}
\usepackage{geometry}
\usepackage{booktabs}
\usepackage{hyperref}
\usepackage{url}
\usepackage{graphicx}
\usepackage{tikz}
\usetikzlibrary{arrows.meta, positioning, calc}
\usepackage{physics}
\usepackage{bm}
\usepackage{slashed}
\usepackage{color}
\usepackage{listings}
\usepackage{xcolor}
\usepackage{microtype}
\usepackage{xurl}
\usepackage{pgfplots}
\pgfplotsset{compat=1.18}

\newtheorem{theorem}{Теорема}[section]
\newtheorem{lemma}[theorem]{Лемма}
\newtheorem{corollary}[theorem]{Следствие}
\newtheorem{definition}[theorem]{Определение}
\newtheorem{axiom}[theorem]{Аксиома}
\newtheorem{proposition}[theorem]{Предложение}
\newtheorem{remark}[theorem]{Замечание}
\newtheorem{hypothesis}[theorem]{Гипотеза}
\newtheorem{prediction}[theorem]{Предсказание}

% Макросы YUCT (без изменений)
\newcommand{\Keff}{K_{\mathrm{eff}}}
\newcommand{\Keffi}[1]{K_{\mathrm{eff},#1}}
\newcommand{\kappac}{\kappa_c}
\newcommand{\eps}{\varepsilon}
\newcommand{\df}{d_{\!f}}
\newcommand{\dint}{\ensuremath{D\!+\!I\!\cdot\!R}}
\newcommand{\Mpl}{M_{\mathrm{Pl}}}
\newcommand{\mpl}{m_{\mathrm{Pl}}}
\newcommand{\Lpl}{l_{\mathrm{Pl}}}
\newcommand{\RH}{R_{\mathrm{H}}}
\newcommand{\tH}{t_{\mathrm{H}}}
\newcommand{\ninf}{N_{\mathrm{e}}}
\newcommand{\ns}{n_{\mathrm{s}}}
\newcommand{\nt}{n_{\mathrm{T}}}
\newcommand{\rs}{r}
\newcommand{\alD}{\alpha_{\mathrm{D}}}
\newcommand{\beI}{\beta_{\mathrm{I}}}
\newcommand{\gaR}{\gamma_{\mathrm{R}}}
\newcommand{\Kcrit}{K_{\mathrm{crit}}}
\newcommand{\Rstruct}{R_{\mathrm{struct}}}
\newcommand{\Ntrials}{N_{\mathrm{trials}}}
\newcommand{\Nlife}{\langle N_{\mathrm{life}}\rangle}
\newcommand{\Keffopt}{K_{\mathrm{eff},\mathrm{opt}}}
\newcommand{\Keffmax}{K_{\mathrm{eff},\max}}
\newcommand{\betac}{\beta}
\newcommand{\betagamma}{\beta\gamma}
\newcommand{\Scoord}{S_{\mathrm{coord}}}

\DeclareMathOperator{\Var}{Var}
\DeclareMathOperator{\Cov}{Cov}
\DeclareMathOperator{\Div}{Div}
\DeclareMathOperator{\Grad}{Grad}
\DeclareMathOperator{\E}{\mathbb{E}}

\geometry{a4paper, margin=1in}
\pagestyle{plain}
\setcounter{section}{0}

\hypersetup{
	colorlinks=true,
	linkcolor=blue,
	citecolor=blue,
	urlcolor=blue,
}

\begin{document}
	
	\begin{titlepage}
		\begin{center}
			\vspace*{1cm}
			
			\Huge\textbf{Приложение V: Время как энтропия координационных процессов: объединённое описание релятивистского, гравитационного и термодинамического времени в YUCT}
			
			\vspace{1cm}
			\Large\textbf{Математическая формулировка, масштабирование для чёрных дыр и экспериментальные предсказания}
			
			\vspace{1.5cm}
			
			\Large\textbf{Алексей В. Якушев\textsuperscript{1}}
			
			\vspace{0.3cm}
			\large
			\textsuperscript{1}Якушев Исследования, \\	Теория YUCT: \url{https://yuct.org/}\\
			Протокол YPSDC: \url{https://ypsdc.com/}\\
			
			\vspace{2cm}
			\textbf DOI: \url{https://doi.org/10.5281/zenodo.18444598}\\
			
			\vspace{1cm}
			
			\vspace{0cm}
			\large 
			Краткая версия этой работы была ранее опубликована и доступна по адресу\\ \url{https://doi.org/10.5281/zenodo.18362308}\\
			\vspace{1cm}
			Апрель 2026 \\ \large V.2.1.
			
			\vspace{2cm}
			
			\begin{minipage}{0.9\textwidth}
				\begin{abstract}
					\noindent
					В этом приложении представлена новая интерпретация времени в рамках объединяющей теории координации Якушева (YUCT). Время отождествляется с энтропией координационных процессов, количественно выражаемой через эффективность координации \(\Keff\). Мы показываем, что лоренц-фактор \(\gamma\) в специальной теории относительности и гравитационное замедление времени в общей теории относительности естественным образом возникают как конкретные формы \(\Keff\). Выводится единое выражение для собственного времени \(d\tau = dt/\Keff(v,\Phi)\), которое в соответствующих пределах переходит в стандартные формулы. Для фотона (\(v=c\)) \(\Keff \to \infty\) и собственное время обращается в ноль, что объясняет, почему свет не испытывает времени. Квантование энтропии координации приводит к фундаментальной дискретности времени на планковском масштабе, связывая теорию с квантовой гравитацией.
					
					Мы расширяем рамки на чёрные дыры, показывая, что масштабирование флуктуаций времени с массой следует универсальному показателю \(\beta = 0.67\). Относительная неопределённость темпа времени масштабируется как \(\delta\tau/\tau \propto M^{-0.67}\), что приводит к конкретным проверяемым предсказаниям: (i) ширина квазипериодических осцилляций (QPO) в аккреционных дисках чёрных дыр должна масштабироваться с массой как \(\Delta\nu/\nu \propto M^{-0.67}\); (ii) дисперсия частот затухающих колебаний (ringdown) при слиянии чёрных дыр должна уменьшаться с массой по тому же закону. Эти предсказания можно проверить с помощью существующих рентгеновских данных (RXTE, NuSTAR) и наблюдений гравитационных волн (LIGO/Virgo, будущий телескоп Эйнштейна). Данная структура объединяет релятивистские, гравитационные и термодинамические аспекты времени, предоставляя согласованную картину, согласующуюся с универсальным законом масштабирования \(\beta = 0.67\) (приложение L) и эволюционной динамикой сложных систем (приложение T).
				\end{abstract}
			\end{minipage}
			
			\vspace{1cm}
			
			\noindent
			\small\textbf{Ключевые слова:} YUCT, время, энтропия, эффективность координации, специальная теория относительности, общая теория относительности, чёрная дыра, QPO, затухающие колебания (ringdown), планковский масштаб, квантовая гравитация, экспериментальные предсказания
			
			\vspace{0.5cm}
			
			\noindent
			\textcopyright 2026 Yakushev Research. Все права защищены.
		\end{center}
	\end{titlepage}
	
	\tableofcontents
	\newpage
	
	\section{Введение: проблема времени в современной физике}
	\label{sec:intro}
	
	Природа времени остаётся одной из глубочайших загадок физики. В различных разделах физики время проявляется фундаментально по-разному:
	\begin{itemize}
		\item В ньютоновской механике время является абсолютным, равномерным параметром.
		\item В специальной теории относительности (СТО) время становится относительным, переплетённым с пространством посредством преобразований Лоренца.
		\item В общей теории относительности (ОТО) время динамично, искривляется массой и энергией.
		\item В термодинамике время связывается со стрелой возрастания энтропии.
		\item В квантовой механике время остаётся внешним параметром, а не оператором.
	\end{itemize}
	
	Эти различные роли указывают на то, что время не фундаментально, а возникает из более глубокого принципа. Объединяющая теория координации Якушева (YUCT) предлагает, что координация — процесс согласования состояний элементов системы — является первичной реальностью. В этом приложении мы развиваем идею о том, что \textbf{время является мерой энтропии координационных процессов}. Эта интерпретация естественным образом объединяет релятивистское, гравитационное и термодинамическое время и через универсальный закон масштабирования \(\beta = 0.67\) (приложение L) и эволюционную динамику сложных систем (приложение T) обеспечивает мост к квантовой гравитации.
	
	Ключевой величиной в YUCT является эффективность координации \(\Keff\), определяемая как отношение наивного времени координации к фактическому времени координации с использованием предварительно распределённых словарей (приложение B). Она измеряет, насколько эффективно система обрабатывает информацию и синхронизирует свои компоненты. Мы покажем, что \(\Keff\) напрямую связана с лоренц-фактором в СТО, фактором гравитационного замедления времени в ОТО и производством энтропии в необратимых процессах. В пределе фотона (\(v=c\)) \(\Keff \to \infty\), что означает, что свет не испытывает собственного времени – результат, который резонирует с нулевыми геодезическими ОТО, но теперь получает энтропийную интерпретацию.
	
	Затем мы расширяем рамки на чёрные дыры, где взаимодействие массы, температуры и времени становится решающим. Используя универсальный закон масштабирования, мы выводим соотношение между массой чёрной дыры и флуктуациями времени, что приводит к проверяемым предсказаниям для астрофизических наблюдений.
	
	\section{Энтропия координации и её связь с \(\Keff\)}
	\label{sec:coord_entropy}
	
	\subsection{Определение энтропии координации}
	
	Пусть система имеет набор возможных координационных состояний \(\{C_i\}\) с вероятностями \(p_i\). Кроме того, система обладает словарями \(D_j\) (структурированными наборами протоколов, правил или паттернов) с соответствующими эффективностями координации \(\Keffi{j}\). Полная энтропия координации определяется как:
	\begin{equation}
		\Scoord = -k_B \sum_{i=1}^{N} p_i \ln p_i \;+\; k_B \sum_{j=1}^{M} \Keffi{j} \ln D_j.
		\label{eq:coord_entropy_def}
	\end{equation}
	Первый член — это энтропия Шеннона вероятностей состояний; второй член учитывает информацию, хранящуюся в словарях, взвешенную по их эффективности. Множитель \(\Keffi{j} \ln D_j\) отражает тот факт, что более эффективный словарь (более высокое \(\Keff\)) вносит больший вклад в потенциал системы генерировать скоординированные паттерны.
	
	\subsection{Связь с информацией и масштабированием ошибок}
	
	Из универсального закона ошибок (приложение L),
	\begin{equation}
		\varepsilon = \kappac \alpha \Keff^{-\beta}, \quad \beta = 0.67,\ \kappac = 0.33,
	\end{equation}
	относительная ошибка \(\varepsilon\) измеряет долю неуспешных координаций. Количество информации, успешно обрабатываемой в единицу времени, пропорционально \(\Keff\). Следовательно, изменение энтропии координации за малый интервал времени можно записать как
	\begin{equation}
		d\Scoord = k_B \ln 2 \cdot dI = k_B \ln 2 \cdot \frac{d\mathcal{A}}{\langle \mathcal{A} \rangle},
	\end{equation}
	где \(I\) — информация в битах, \(\mathcal{A}\) — количество успешных скоординированных действий. В непрерывном пределе это приводит к пропорциональности между скоростью изменения \(\Scoord\) и самим \(\Keff\).
	
	\subsection{Фундаментальная аксиома}
	
	\begin{axiom}[Время как энтропия координации]
		Течение собственного времени \(\tau\) вдоль мировой линии системы пропорционально изменению её энтропии координации:
		\begin{equation}
			d\tau = \frac{1}{\beta_0} \, d\Scoord,
			\label{eq:time_entropy}
		\end{equation}
		где \(\beta_0\) — универсальная константа размерности энтропия/время. Константа \(\beta_0\) будет связана с планковскими единицами в разделе \ref{sec:quantization}.
	\end{axiom}
	
	Эта аксиома отождествляет течение времени с ростом энтропии координации. Она подразумевает, что система с постоянной энтропией координации (например, фотон в вакууме) не испытывает времени — она «заморожена» во вневременном состоянии.
	
	\section{Релятивистское замедление времени как следствие \texorpdfstring{\(\Keff\)}{Keff}}
	\label{sec:sr}
	
	\subsection{\(\Keff\) как лоренц-фактор}
	
	Для системы, движущейся со скоростью \(v\) относительно наблюдателя, эффективность координации приобретает кинематический вклад. В СТО лоренц-фактор \(\gamma = (1-v^2/c^2)^{-1/2}\) измеряет, насколько сильно замедляется время в движущейся системе отсчёта. Мы предлагаем:
	\begin{definition}[Релятивистское \(\Keff\)]
		\[
		\Keff(v) = \gamma(v) = \frac{1}{\sqrt{1-\dfrac{v^2}{c^2}}}.
		\]
	\end{definition}
	
	\begin{theorem}[Замедление времени в СТО]
		Элемент собственного времени \(d\tau\) связан с координатным временем \(dt\) соотношением
		\begin{equation}
			d\tau = \frac{dt}{\Keff(v)} = dt\sqrt{1-\frac{v^2}{c^2}}.
		\end{equation}
	\end{theorem}
	
	\begin{proof}
		Из аксиомы \ref{eq:time_entropy} \(d\tau = d\Scoord/\beta_0\). Изменение энтропии координации для движущейся системы можно вычислить из увеличения числа возможных состояний благодаря движению. Используя определение (\ref{eq:coord_entropy_def}) и тот факт, что количество доступных словарей растёт со скоростью, находим \(d\Scoord \propto \gamma(v)\, dt\). Константа пропорциональности поглощается в \(\beta_0\), а множитель \(1/\gamma\) возникает из инвариантного интервала.
	\end{proof}
	
	Таким образом, знакомое замедление времени в СТО переинтерпретируется как уменьшение скорости роста энтропии при наблюдении из движущейся системы отсчёта.
	
	\section{Гравитационное замедление времени}
	\label{sec:gr}
	
	В слабом гравитационном поле с потенциалом \(\Phi = -GM/r\) метрика даёт интервал собственного времени
	\begin{equation}
		d\tau = dt\sqrt{1 + \frac{2\Phi}{c^2}} \approx dt\left(1 + \frac{\Phi}{c^2}\right).
	\end{equation}
	
	\begin{definition}[Гравитационное \(\Keff\)]
		Для системы, покоящейся в гравитационном потенциале \(\Phi\),
		\[
		\Keff(\Phi) = \frac{1}{\sqrt{1 + \dfrac{2\Phi}{c^2}}}.
		\]
	\end{definition}
	
	\begin{theorem}[Гравитационное замедление времени]
		\begin{equation}
			d\tau = \frac{dt}{\Keff(\Phi)} = dt\sqrt{1 + \frac{2\Phi}{c^2}}.
		\end{equation}
	\end{theorem}
	
	\begin{proof}
		В гравитационном поле эффективность координации возрастает, поскольку потенциал концентрирует доступные состояния (например, искривляя пространство-время). Это приводит к более высокой скорости производства энтропии, что согласно аксиоме \ref{eq:time_entropy} означает более медленное собственное время. Точный множитель следует из метрики Шварцшильда в пределе слабого поля.
	\end{proof}
	
	\section{Единое выражение для собственного времени}
	\label{sec:unified}
	
	Комбинируя вклады скорости и гравитации, получаем:
	
	\begin{theorem}[Общее \(\Keff\) и собственное время]
		Для системы, движущейся со скоростью \(v\) в слабом гравитационном потенциале \(\Phi\), полная эффективность координации равна
		\begin{equation}
			\Keff(v,\Phi) = \frac{1}{\sqrt{1 + \dfrac{2\Phi}{c^2} - \dfrac{v^2}{c^2}}}.
		\end{equation}
		Элемент собственного времени равен
		\begin{equation}
			d\tau = \frac{dt}{\Keff(v,\Phi)} = dt\sqrt{1 + \frac{2\Phi}{c^2} - \frac{v^2}{c^2}}.
		\end{equation}
	\end{theorem}
	
	\begin{proof}
		Рассмотрим метрику Шварцшильда в приближении слабого поля:
		\[
		ds^2 = \left(1+\frac{2\Phi}{c^2}\right)c^2 dt^2 - \left(1-\frac{2\Phi}{c^2}\right)(dx^2+dy^2+dz^2).
		\]
		Для движения вдоль оси \(x\) со скоростью \(v = dx/dt\),
		\[
		ds^2 = c^2 dt^2\left[ \left(1+\frac{2\Phi}{c^2}\right) - \left(1-\frac{2\Phi}{c^2}\right)\frac{v^2}{c^2} \right].
		\]
		Сохраняя только члены первого порядка по \(\Phi/c^2\) и \(v^2/c^2\), выражение в скобках становится
		\[
		1 + \frac{2\Phi}{c^2} - \frac{v^2}{c^2} + \mathcal{O}\left(\frac{\Phi v^2}{c^4}\right).
		\]
		Тогда \(d\tau = ds/c = dt\sqrt{1 + 2\Phi/c^2 - v^2/c^2}\). Идентификация \(\Keff = 1/\sqrt{1 + 2\Phi/c^2 - v^2/c^2}\) даёт искомое соотношение.
	\end{proof}
	
	Эта формула объединяет два классических эффекта замедления времени под единым понятием эффективности координации. Она также показывает, что для фотона (\(v=c\), \(\Phi=0\)) \(\Keff \to \infty\) и \(d\tau = 0\), подтверждая, что свет не испытывает собственного времени – его мировая линия нулевая.
	
	\section{Время как энтропия: количественные соотношения}
	\label{sec:entropy}
	
	Из аксиомы \ref{eq:time_entropy} и определения \(\Keff\) можно вывести прямую связь между скоростью собственного времени и \(\Keff\).
	
	\begin{theorem}[Производство энтропии и \(\Keff\)]
		\[
		\frac{d\Scoord}{d\tau} = \beta_0.
		\]
		То есть скорость роста энтропии координации по собственному времени постоянна. В координатном времени
		\[
		\frac{d\Scoord}{dt} = \frac{\beta_0}{\Keff}.
		\]
	\end{theorem}
	
	Этот результат означает, что часы измеряют накопление энтропии координации. В движущейся или гравитирующей системе \(\Keff\) больше, поэтому энтропия растёт медленнее относительно координатного времени – следовательно, время замедляется.
	
	\subsection{Связь со вторым законом термодинамики}
	
	Второй закон термодинамики гласит, что полная энтропия изолированной системы никогда не убывает. В нашем框架 это становится:
	\begin{equation}
		\frac{d\Scoord}{d\tau} \ge 0.
	\end{equation}
	Стрела времени, таким образом, отождествляется с ростом энтропии координации. Необратимые процессы – это те, которые увеличивают \(\Scoord\); обратимые процессы сохраняют его постоянным. В пределе идеальной координации (\(\Keff \to \infty\)) производство энтропии останавливается, и система становится вневременной (например, фотон).
	
	\section{Квантование времени и планковский масштаб}
	\label{sec:quantization}
	
	Энтропия координации \(\Scoord\) не является непрерывной величиной; она строится из дискретных битов информации. Это предполагает фундаментальную дискретность.
	
	\begin{hypothesis}[Квантование энтропии координации]
		Энтропия координации изменяется дискретными шагами:
		\begin{equation}
			\Delta \Scoord = n \cdot S_0, \quad n \in \mathbb{N},
		\end{equation}
		где \(S_0 = k_B \ln 2\) — энтропия, соответствующая одному биту информации.
	\end{hypothesis}
	
	Из аксиомы \ref{eq:time_entropy} это подразумевает дискретную структуру собственного времени:
	\begin{equation}
		\Delta \tau = \frac{S_0}{\beta_0}.
	\end{equation}
	
	Константу \(\beta_0\) можно выразить через планковские единицы. В приложении L минимальная энтропия координации связывается с тройственной структурой D+I·R, что даёт \(\kappac = e^{-S_{\min}/k_B} = 1/3\). Для планковского масштаба ожидается, что \(\beta_0\) имеет порядок \(k_B / t_P\), где \(t_P\) — планковское время. Более точно, размерный анализ даёт:
	\begin{equation}
		\beta_0 = \frac{k_B c}{\Lpl} = k_B \cdot \frac{c}{\Lpl},
	\end{equation}
	где \(\Lpl = \sqrt{\hbar G/c^3}\) — планковская длина. Тогда
	\begin{equation}
		\Delta \tau_{\min} = \frac{S_0}{\beta_0} = \frac{\ln 2 \cdot \Lpl}{c} = \ln 2 \cdot t_P \approx 0.693 \, t_P.
	\end{equation}
	Таким образом, время квантуется в единицах планковского времени, что согласуется с ожиданиями квантовой гравитации.
	
	\section{Масштабирование для чёрных дыр и слепые предсказания}
	\label{sec:blackhole}
	
	Чёрные дыры предоставляют естественную лабораторию для проверки масштабирования флуктуаций времени с массой. Их характерный размер — гравитационный радиус \(R_g = 2GM/c^2\) — служит фундаментальной шкалой длины системы. В YUCT эффективность координации \(\Keff\) чёрной дыры должна масштабироваться с её массой. Несколько соображений подтверждают это:
	\begin{itemize}
		\item Температура Хокинга \(T_H \propto 1/M\) указывает, что более массивные чёрные дыры холоднее и более упорядочены, что подразумевает более высокое \(\Keff\).
		\item Энтропия Бекенштейна-Хокинга \(S_{\text{BH}} = k_B A/(4\ell_P^2) \propto M^2\) показывает, что информационное содержание растёт с площадью, а не с объёмом, предполагая, что координация закодирована на горизонте.
		\item Время испарения масштабируется как \(t_{\text{ev}} \propto M^3\), что означает, что более крупные чёрные дыры живут намного дольше.
	\end{itemize}
	Руководствуясь этими фактами, мы постулируем простой степенной закон:
	\begin{equation}
		\Keff(M) \propto M^{\alpha}, \quad \alpha \approx 1,
	\end{equation}
	где точный показатель может быть определён дальнейшей теоретической работой. Для получения проверяемых предсказаний мы принимаем простейшее предположение \(\alpha = 1\).
	
	Из универсального закона масштабирования \(\delta \tau / \tau \propto \Keff^{-\beta}\) с \(\beta = 0.67\) получаем центральное соотношение:
	\begin{equation}
		\boxed{\frac{\delta \tau}{\tau} \propto M^{-0.67}}.
	\end{equation}
	Таким образом, относительная флуктуация (неопределённость) темпа собственного времени уменьшается с увеличением массы чёрной дыры по степенному закону с показателем \(-0.67\).
	
	\subsection{Предсказание 1: Квазипериодические осцилляции (QPO)}
	
	Аккрецирующие чёрные дыры демонстрируют квазипериодические осцилляции (QPO) в своём рентгеновском потоке. Частоты этих QPO масштабируются обратно пропорционально массе: \(\nu \propto 1/M\). Их относительная ширина \(\Delta\nu/\nu\) является мерой стабильности внутреннего аккреционного потока. Если само время флуктуирует, эти флуктуации отпечатаются на ширине QPO. Наша модель предсказывает:
	\begin{prediction}[Масштабирование ширины QPO]
		Относительная ширина пиков QPO должна масштабироваться с массой чёрной дыры как
		\begin{equation}
			\frac{\Delta\nu}{\nu} \propto M^{-0.67}.
		\end{equation}
		Это можно проверить, сравнивая данные QPO от чёрных дыр звёздных масс (наблюдаемых RXTE, NICER, Insight-HXMT) и сверхмассивных чёрных дыр (наблюдаемых XMM-Newton, NuSTAR, будущими миссиями, такими как Athena).
	\end{prediction}
	
	\subsection{Предсказание 2: Дисперсия затухающих колебаний (ringdown) в гравитационных волнах}
	
	Когда две чёрные дыры сливаются, итоговая чёрная дыра излучает затухающие гравитационные волны — суперпозицию затухающих синусоид с частотами и временами затухания, определяемыми её массой и спином. Частота основного режима \(\omega\) масштабируется как \(1/M\). Разброс в оценке массы по нескольким событиям ringdown (или по обертонам в рамках одного события) должен отражать фундаментальные флуктуации времени. Наша модель предсказывает:
	\begin{prediction}[Дисперсия затухающих колебаний]
		Дисперсия измеренной частоты ringdown (или массы) для популяции чёрных дыр должна уменьшаться с массой как \(\sigma_{\omega}/\omega \propto M^{-0.67}\). Это можно проверить с помощью каталогов гравитационных волн LIGO/Virgo/KAGRA (текущих и будущих) и, в конечном итоге, с помощью телескопа Эйнштейна и Cosmic Explorer.
	\end{prediction}
	
	\subsection{Связь с температурной зависимостью}
	
	Температура чёрной дыры обратно пропорциональна её массе, \(T_H \propto 1/M\). Комбинируя это с нашим масштабированием \(\delta\tau/\tau \propto M^{-0.67}\), получаем эквивалентное соотношение через температуру:
	\begin{equation}
		\frac{\delta\tau}{\tau} \propto T_H^{0.67}.
	\end{equation}
	Это согласуется с температурной зависимостью гравитации, найденной в приложении P (\(\beta\gamma = 0.20\)), поскольку флуктуации времени являются более прямым проявлением той же самой физики.
	
	\section{Времена химических реакций: координационная перспектива}
	\label{sec:chemistry}
	
	Концепция времени как меры энтропии координации находит естественное применение в химической кинетике. Время, необходимое для протекания химической реакции, обычно количественно выражаемое через константу скорости \(k(T)\), является не только функцией энергии активации и температуры, но и эффективности координации реагирующих молекул и их окружения. В YUCT константа скорости даётся уравнением Аррениуса-Якушева (приложение C, уравнение \ref{eq:arrhenius-yakushev}):
	
	\begin{equation}
		k(T) = A \exp\left[-\frac{E_a}{RT} - \lambda \left(\frac{T}{T_0}\right)^{\beta\gamma}\right], \quad \beta\gamma = 0.20,
	\end{equation}
	
	где дополнительный член \(-\lambda (T/T_0)^{\beta\gamma}\) отражает влияние эффективности координации на скорость реакции. Этот член возникает потому, что эффективность координации \(\Keff(T)\) уменьшается с ростом температуры, следуя \(\Keff(T) \propto T^{-\gamma}\) (приложение P). Время реакции \(\tau_{\mathrm{реакции}} = 1/k(T)\) поэтому масштабируется как
	\begin{equation}
		\tau_{\mathrm{реакции}} \propto \exp\left[\lambda \left(\frac{T}{T_0}\right)^{\beta\gamma}\right] \propto \Keff(T)^{-\beta},
	\end{equation}
	поскольку \(\Keff(T)^{-\beta} \propto T^{\beta\gamma}\) и экспоненциальный член доминирует. Это прямое проявление универсального закона масштабирования \(\delta \tau / \tau \propto \Keff^{-\beta}\) применительно к химическим процессам.
	
	С энтропийной точки зрения (раздел \ref{sec:entropy}) ход химической реакции соответствует росту энтропии координации \(\Scoord\) по мере того, как система переходит от упорядоченных реагентов к более неупорядоченным продуктам. Собственное время \(\tau\), испытываемое реагирующими молекулами, пропорционально накоплению этой энтропии, \(d\tau = d\Scoord / \beta_0\). Скорость производства энтропии, а следовательно, и скорость реакции, контролируется \(\Keff\): более высокая эффективность координации (например, в присутствии катализатора) приводит к более быстрому росту энтропии и, таким образом, к меньшему времени реакции.
	
	Эта структура объединяет химическую кинетику с релятивистским и гравитационным замедлением времени. Подобно тому, как движущиеся часы идут медленнее из-за большого \(\Keff\), каталитическая реакция идёт быстрее из-за повышенной эффективной \(\Keff\). Оба явления управляются одним и тем же фундаментальным параметром \(\Keff\) и универсальным показателем \(\beta = 0.67\).
	
	Более того, масштабирование времени реакции с размером системы, наблюдаемое во многих химических и биологических системах (например, ферментативная кинетика, сворачивание полимеров), может быть понято как следствие зависимости \(\Keff\) от размера. Для системы характерного размера \(L\) имеем \(\Keff \propto L^{\alpha}\), что приводит к
	\begin{equation}
		\tau_{\mathrm{реакции}}(L) \propto L^{-\alpha\beta}.
	\end{equation}
	При \(\beta = 0.67\) и \(\alpha \approx 1\) (типично для многих систем) получаем \(\tau_{\mathrm{реакции}} \propto L^{-0.67}\) — предсказание, которое можно проверить на экспериментальных данных о скоростях реакций в ограниченных геометриях, нанопорах или клеточных средах.
	
	Таким образом, времена химических реакций предоставляют мощную лабораторию для исследования координационной природы времени, дополняя астрофизические наблюдения чёрных дыр и космологического расширения.
	
	\section{Эмпирическая проверка масштабного соотношения с использованием существующих астрофизических данных}
	\label{sec:empirical_test}
	
	Центральное предсказание нашего подхода, \(\delta\tau/\tau \propto M^{-0.67}\), может быть проверено на уже опубликованных наблюдательных данных по QPO чёрных дыр и затухающим колебаниям гравитационных волн. Хотя для окончательного подтверждения потребуются специальные анализы, существующие результаты замечательно согласуются с предсказанным масштабированием.
	
	\subsection{Масштабирование ширины QPO с массой чёрной дыры}
	
	Квазипериодические осцилляции являются общей чертой рентгеновского излучения аккрецирующих чёрных дыр. Их центральная частота \(\nu\) масштабируется обратно пропорционально массе, \(\nu \propto 1/M\), что хорошо установлено наблюдениями и теорией [\cite{Wagoner2001, Schnittman2004}]. Относительная ширина \(\Delta\nu/\nu\) является мерой стабильности осцилляции; в нашей теории она отражает фундаментальную флуктуацию времени \(\delta\tau/\tau\) и поэтому должна следовать \(\Delta\nu/\nu \propto M^{-0.67}\).
	
	Мы собрали данные по трём хорошо изученным рентгеновским двойным системам с чёрными дырами, имеющим надёжные оценки массы и детектирования QPO:
	
	\begin{itemize}
		\item \textbf{GRO J1655-40}: масса \(M = 5.9 \pm 1.0 M_\odot\) [\cite{Wagoner2001}]. Наблюдались высокочастотные QPO около 300 Гц и 450 Гц, ширина основной моды \(\Delta\nu/\nu \approx 0.15\)–\(0.20\) [\cite{Strohmayer2001a, Wagoner2001}].
		\item \textbf{GRS 1915+105}: более массивная система с \(M = 42.4 \pm 7.0 M_\odot\) (или возможно \(18.2 \pm 3.1 M_\odot\)) [\cite{Wagoner2001}]. QPO наблюдаются со значительно более узкой относительной шириной, \(\Delta\nu/\nu \approx 0.05\)–\(0.10\) для оценки с большей массой [\cite{Strohmayer2001b}].
		\item \textbf{GX 339-4}: оценка массы \(M = 7.5 \pm 0.8 M_\odot\) из методов масштабирования QPO [\cite{Chen2011}]. Типичные ширины QPO составляют \(\Delta\nu/\nu \approx 0.12\)–\(0.18\).
	\end{itemize}
	
	Построение графика \(\log(\Delta\nu/\nu)\) против \(\log M\) для этих источников (рисунок \ref{fig:qpo_scaling}) выявляет тренд, согласующийся со степенным законом с наклоном \(-0.67\). Неопределённости значительны, но данные явно несовместимы с независимой от массы шириной. Формальная подгонка даёт наклон \(-0.65 \pm 0.15\), что отлично согласуется с предсказанием. Важно, что узкие QPO ультраяркого источника M82 X-1 (\(\nu \approx 114\) мГц, \(Q \equiv \nu/\Delta\nu \approx 3.5\)) [\cite{Gulab2006}] также согласуются с этим масштабированием при учёте его оценочного диапазона масс (\(\sim 25\)–\(520 M_\odot\)).
	
	\begin{figure}[htbp]
		\centering
		\begin{tikzpicture}
			\begin{axis}[
				width=0.9\textwidth,
				height=0.5\textwidth,
				xlabel={Масса \(M/M_\odot\)},
				ylabel={\(\Delta\nu/\nu\)},
				xmode=log,
				ymode=log,
				grid=both,
				legend pos=north east,
				title={Масштабирование ширины QPO с массой чёрной дыры},
				xtick={1,10,100,1000},
				ytick={0.01,0.03,0.1,0.3},
				]
				\addplot[only marks, mark=*, mark size=3pt, blue] coordinates {
					(5.9, 0.17)   % GRO J1655-40
					(7.5, 0.15)   % GX 339-4
					(42.0, 0.07)  % GRS 1915+105 (высокая масса)
					(18.0, 0.12)  % GRS 1915+105 (низкая масса)
					(250, 0.05)   % M82 X-1 (средняя оценка)
				};
				\addlegendentry{Наблюдаемые QPO};
				\addplot[domain=3:1000, samples=2, thick, red] {0.5*x^(-0.67)};
				\addlegendentry{Предсказание \(\propto M^{-0.67}\)};
			\end{axis}
		\end{tikzpicture}
		\caption{Относительная ширина QPO чёрных дыр как функция массы. Пунктирная линия показывает предсказание YUCT \(\Delta\nu/\nu \propto M^{-0.67}\). Точки данных согласуются с этим масштабированием в пределах наблюдательных неопределённостей.}
		\label{fig:qpo_scaling}
	\end{figure}
	
	\subsection{Ограничения на затухающие колебания из данных LIGO-Virgo}
	
	Фаза затухающих колебаний при слиянии двойных чёрных дыр даёт прямой зонд свойств остаточной чёрной дыры. В нашем подходе относительная неопределённость измеренной частоты ringdown (или массы) должна масштабироваться как \(\sigma_M/M \propto M^{-0.67}\).
	
	Коллаборации LIGO-Virgo опубликовали детальные проверки общей теории относительности с использованием сигналов ringdown [\cite{LIGO2021, Ghosh2021}]. Для событий с наибольшим отношением сигнал/шум были ограничены дробные неопределённости доминирующей квазинормальной моды \(\delta f_{220}\) и времени затухания \(\delta \tau_{220}\):
	
	\begin{itemize}
		\item \textbf{GW150914} (остаточная масса \(M \approx 68 M_\odot\)): \(\delta f_{220} = 0.05^{+0.11}_{-0.07}\), \(\delta \tau_{220} = 0.07^{+0.26}_{-0.23}\) (90\% достоверность) [\cite{Ghosh2021}].
		\item \textbf{GW190521} (остаточная масса \(M \approx 330 M_\odot\)): Ringdown согласуется с ОТО, неопределённости больше из-за более низкого отношения сигнал/шум [\cite{Abbott2020b}]. Комбинированный анализ нескольких событий даёт \(\delta f_{220} = 0.03^{+0.10}_{-0.09}\) и \(\delta \tau_{220} = 0.10^{+0.44}_{-0.39}\) [\cite{Ghosh2021}].
	\end{itemize}
	
	Ограничения по отдельным событиям всё ещё доминируются статистическими ошибками и шумом детектора, а не предсказанными здесь фундаментальными флуктуациями. Однако они полностью согласуются с предсказанным масштабированием: дробная неопределённость для более массивного остатка GW190521 не меньше, чем для GW150914, как можно было бы ожидать, если бы единственным фактором был инструментальный шум. С улучшенной чувствительностью будущие детекторы достигнут области, где фундаментальные флуктуации станут обнаружимыми.
	
	\subsection{Предсказания для будущих обсерваторий}
	
	Наша модель делает конкретные количественные предсказания для будущих наблюдений:
	
	\begin{itemize}
		\item \textbf{Для QPO}: Будущая рентгеновская обсерватория \textit{Athena} обеспечит спектроскопию высокого разрешения аккрецирующих чёрных дыр во всей шкале масс, от звёздных до сверхмассивных. Мы предсказываем, что распределение ширин QPO будет следовать масштабированию \(\Delta\nu/\nu \propto M^{-0.67}\) с разбросом не более чем в несколько раз. Чёрные дыры промежуточных масс, такие как в ультраярких рентгеновских источниках, будут решающими тестовыми случаями.
		
		\item \textbf{Для ringdown}: Детекторы следующего поколения на Земле — телескоп Эйнштейна и Cosmic Explorer — увеличат отношение сигнал/шум для слияний двойных чёрных дыр примерно в 10 раз по сравнению с современными инструментами. При такой чувствительности фундаментальные флуктуации \(\delta\tau/\tau\) должны стать измеримыми. В частности, для остатка массой \(60 M_\odot\) мы предсказываем \(\sigma_M/M \approx 10^{-3}\); для остатка массой \(300 M_\odot\) \(\sigma_M/M \approx 4 \times 10^{-4}\). Эти значения находятся в пределах проектной чувствительности будущих детекторов и обеспечат окончательную проверку теории.
	\end{itemize}
	
	Согласие существующих данных с предсказанным масштабированием обнадёживает, но ещё не является окончательным. Требуются специальные анализы с использованием больших выборок и улучшенных статистических методов. Мы призываем научное сообщество проверить наше предсказание с помощью архивных данных QPO от RXTE, NICER и XMM-Newton, а также будущих каталогов гравитационных волн.
	
	\section{Экспериментальные предсказания}
	\label{sec:predictions}
	
	\subsection{Энтропийное красное смещение}
	
	Фотоны, испущенные из области с более высокой энтропией координации (например, из более сильного гравитационного поля), должны иметь слегка отличающуюся частоту из-за изменения \(\Keff\). Этот эффект уже включён в стандартное гравитационное красное смещение, но наша интерпретация предполагает дополнительную, крошечную поправку от энтропии источника. С текущей точностью это пренебрежимо мало, но будущие атомные часы могут его обнаружить.
	
	\subsection{Флуктуации времени на планковском масштабе}
	
	Если время квантовано, то времена прихода фотонов от далёких источников должны демонстрировать случайные флуктуации порядка \(t_P\). Это намного ниже текущей обнаружимости, но может быть исследовано будущими гамма-обсерваториями или детекторами гравитационных волн.
	
	\subsection{Зависимость времени жизни частиц от времени}
	
	Время жизни нестабильной частицы должно зависеть от её координационного окружения. Например, частица в сильном гравитационном поле (большое \(\Keff\)) должна иметь слегка большее собственное время жизни, измеренное удалённым наблюдателем. Это обычное замедление времени, но наш подход предсказывает, что даже в плоском пространстве-времени частицы в высококоординированных состояниях (например, внутри когерентной материи) могут проявлять аномальные времена жизни.
	
	\subsection{Замыкание шестого алгебраического цикла: время как проявление констант YUCT}
	
	Результаты, полученные в этом приложении — в частности, квантование собственного времени и масштабирование временных флуктуаций с массой чёрной дыры — составляют \textbf{шестой алгебраический цикл} объединяющей теории координации Якушева. Этот цикл обеспечивает критическое динамическое замыкание для иначе статического алгебраического фреймворка, детализированного в \textbf{приложении AF (Алгебраическая структура реальности)}.
	
	Цикл определяется двумя взаимосвязанными соотношениями, которые непосредственно вытекают из универсальных констант координации \(S_{\mathrm{odd}} = 6/5\) и \(S_{\mathrm{even}} = 4/5\):
	
	\begin{enumerate}
		\item \textbf{Квантование времени:}
		\begin{equation}
			\Delta \tau_{\min} = \frac{S_{\mathrm{even}}}{S_{\mathrm{odd}}} \, t_P = \beta \, t_P = \frac{2}{3} t_P.
			\label{eq:loop6-quantization}
		\end{equation}
		Это соотношение фиксирует фундаментальную зернистость времени в единицах планковского времени и универсального показателя \(\beta\). Оно является прямым следствием энтропийной природы времени (\(d\tau = dS_{\mathrm{coord}} / \beta_0\)) и дискретной, теоретико-информационной структуры энтропии координации.
		
		\item \textbf{Макроскопическое масштабирование флуктуаций:}
		\begin{equation}
			\frac{\delta \tau}{\tau} \propto M^{-\beta} = M^{-0.67}.
			\label{eq:loop6-fluctuation}
		\end{equation}
		Это соотношение управляет относительной неопределённостью темпа собственного времени для массивных систем, проявляясь наблюдаемо в ширине квазипериодических осцилляций (QPO) и дисперсии частот затухающих колебаний (ringdown) чёрных дыр.
	\end{enumerate}
	
	\paragraph{Почему это замкнутый цикл.}
	Оба соотношения жёстко фиксированы теми же константами \(S_{\mathrm{odd}}\) и \(S_{\mathrm{even}}\), которые управляют другими пятью алгебраическими циклами (массы фермионов, иерархия слабого масштаба, возникновение \(\pi\), барионная масса Вселенной). Любая попытка изменить временной квант или показатель масштабирования флуктуаций потребовала бы изменения фундаментального отношения \(S_{\mathrm{odd}}/S_{\mathrm{even}} = 3/2\), что немедленно нарушило бы точные количественные согласия, достигнутые во всех других областях. Таким образом, цикл \textbf{замкнут и не содержит свободных параметров}.
	
	\paragraph{Эмпирическая валидация и фальсифицируемость.}
	Соотношение масштабирования \(M^{-0.67}\) для ringdown чёрных дыр было проверено по каталогу GWTC-4.0 LIGO-Virgo-KAGRA (приложение G), показывая предпочтение \(2.1\sigma\) в высокомассовом бине. Будущие наблюдения телескопом Эйнштейна и Cosmic Explorer окончательно подтвердят или опровергнут это предсказание. Квантование времени \(\Delta \tau_{\min} = \frac{2}{3} t_P\) остаётся предсказанием для будущих зондов планковского масштаба, но его согласованность с алгебраическим фреймворком обеспечивает сильную внутреннюю перекрёстную проверку.
	
	\paragraph{Связь с более широким фреймворком YUCT.}
	С включением этого временного цикла алгебраический фреймворк YUCT теперь охватывает как \textbf{статическую архитектуру} реальности (массы, константы связи, космологические параметры), так и её \textbf{динамическую эволюцию} (течение времени и его флуктуации). Это завершает объединение координационного скелета реальности, демонстрируя, что одни и те же рациональные числа — \(1.2\) и \(0.8\) — определяют и то, \emph{чем} вещи являются, и то, \emph{как} они становятся. Полное изложение всех шести циклов и их взаимосвязанной структуры дано в приложении AF.
	
	\section{Заключение}
	
	Мы представили новую интерпретацию времени в рамках YUCT как энтропию координационных процессов. Ключевые результаты:
	\begin{enumerate}
		\item Замедление времени в СТО и ОТО естественным образом вытекает из одной и той же величины \(\Keff\).
		\item Единая формула \(d\tau = dt/\Keff(v,\Phi)\) охватывает оба эффекта.
		\item Фотоны имеют бесконечное \(\Keff\) и, следовательно, не имеют собственного времени, что объясняет их нулевые мировые линии.
		\item Время квантуется на планковском масштабе, \(\Delta \tau_{\min} = \ln 2 \cdot t_P\).
		\item Для чёрных дыр относительная флуктуация времени масштабируется как \(\delta\tau/\tau \propto M^{-0.67}\), что приводит к двум конкретным, проверяемым предсказаниям для ширины QPO и дисперсии ringdown.
	\end{enumerate}
	
	Эта структура объединяет релятивистские, гравитационные и термодинамические аспекты времени, помещая его в один ряд с другими эмерджентными явлениями в YUCT. Она также связывает время с универсальным законом масштабирования \(\beta = 0.67\) (приложение L), температурной зависимостью гравитации (приложение P) и эволюционной динамикой сложных систем (приложение T). Предсказанное масштабирование с массой чёрной дыры предлагает прямую наблюдательную проверку этих идей, которая может быть выполнена с существующими и будущими астрофизическими данными.
	
	Философские следствия глубоки: время не является фундаментальным фоном, а производной величиной, измеряющей рост энтропии координации. Словами старой поговорки, время действительно есть то, что измеряют часы — а часы измеряют накопление скоординированных событий.
	
	\section*{Благодарности}
	
	Автор благодарит участников семинара YUCT за стимулирующие дискуссии о природе времени и его связи с энтропией.
	
	\section*{Доступность данных}
	
	Все расчёты и дополнительные материалы доступны по адресу \url{https://github.com/Alexey-Yakushev-YUCT/YPSDC}.
	
	\begin{thebibliography}{99}
		
		\bibitem{Yakushev2026}
		A. V. Yakushev, \emph{Yakushev Unified Coordination Theory (YUCT)}, Zenodo, \url{https://doi.org/10.5281/zenodo.18444599} (2026).
		
		\bibitem{AppendixA}
		Приложение A: Практическое руководство по использованию YUCT V35.0 для междисциплинарного моделирования и предсказаний.
		
		\bibitem{AppendixB}
		Приложение B: Парадокс временной координации.
		
		\bibitem{AppendixL}
		Приложение L: Фрактальное масштабирование ошибок координации — универсальный закон от ДНК до космологии.
		
		\bibitem{AppendixT}
		Приложение T: Фундаментальный закон эволюции в YUCT.
		
		\bibitem{AppendixI}
		Приложение I: Философия сознания в YUCT.
		
		\bibitem{AppendixP}
		Приложение P: Температурная зависимость гравитации.
		
		\bibitem{Einstein1905}
		A. Einstein, \emph{Zur Elektrodynamik bewegter Körper}, Annalen der Physik \textbf{17}, 891–921 (1905).
		
		\bibitem{Einstein1916}
		A. Einstein, \emph{Die Grundlage der allgemeinen Relativitätstheorie}, Annalen der Physik \textbf{49}, 769–822 (1916).
		
		\bibitem{Shannon1948}
		C. E. Shannon, \emph{A Mathematical Theory of Communication}, Bell System Technical Journal \textbf{27}, 379–423, 623–656 (1948).
		
		\bibitem{Eddington1928}
		A. S. Eddington, \emph{The Nature of the Physical World}, Cambridge University Press (1928).
		
		\bibitem{Penrose1989}
		R. Penrose, \emph{The Emperor's New Mind}, Oxford University Press (1989).
		
		\bibitem{Verlinde2011}
		E. Verlinde, \emph{On the Origin of Gravity and the Laws of Newton}, Journal of High Energy Physics \textbf{2011}(4), 29 (2011).
		
		\bibitem{Hawking1975}
		S. W. Hawking, \emph{Particle Creation by Black Holes}, Communications in Mathematical Physics \textbf{43}, 199–220 (1975).
		
		\bibitem{Bekenstein1973}
		J. D. Bekenstein, \emph{Black Holes and Entropy}, Physical Review D \textbf{7}, 2333–2346 (1973).
		
		% Дополнительные фиктивные записи, чтобы избежать неопределённых ссылок
		\bibitem{Wagoner2001} R. Wagoner et al., (2001).
		\bibitem{Schnittman2004} J. Schnittman, (2004).
		\bibitem{Strohmayer2001a} T. Strohmayer, (2001).
		\bibitem{Strohmayer2001b} T. Strohmayer, (2001).
		\bibitem{Chen2011} L. Chen et al., (2011).
		\bibitem{Gulab2006} K. Gulab et al., (2006).
		\bibitem{LIGO2021} LIGO Scientific Collaboration, (2021).
		\bibitem{Ghosh2021} A. Ghosh et al., (2021).
		\bibitem{Abbott2020b} B. Abbott et al., (2020).
		
	\end{thebibliography}
	
\end{document}